\documentclass{ximera}

\input{../../../preamble.tex}


\definecolor{fvdnavy}{HTML}{071D43}
\definecolor{fvdcyan}{HTML}{20BFC6}
\definecolor{fvdcyanlight}{HTML}{EFFCFB}
\definecolor{fvdmagenta}{HTML}{F10D69}
\definecolor{fvdmagentalight}{HTML}{FFF1F7}
\definecolor{fvdtext}{HTML}{163044}





%=================================================
% Pedagogische omgevingen
% PDF: tcolorbox
% HTML: learning-box
%=================================================

\newenvironment{voorkennis}
{%
    \ifdefined\HCode
        \HCode{%
<section class="learning-box learning-box--prior">
<div class="learning-box__header">
<span class="learning-box__icon" aria-hidden="true"></span>
<span class="learning-box__title">Voorkennis</span>
</div>
<div class="learning-box__content">
        }%
        \begin{itemize}
    \else
       \begin{tcolorbox}[
    enhanced,
    breakable,
    colback=fvdcyanlight,
    colframe=fvdcyan,
    coltext=fvdtext,
    boxrule=0.5pt,
    leftrule=4pt,
    arc=4mm,
    outer arc=4mm,
    left=7mm,
    right=6mm,
    top=5mm,
    bottom=5mm,
    before skip=6mm,
    after skip=6mm,
    title=Voorkennis,
    coltitle=fvdcyan,
    fonttitle=\bfseries\Large,
    attach title to upper={\par\smallskip},
    boxed title style={
        empty,
        left=0mm,
        right=0mm,
        top=0mm,
        bottom=0mm
    },
    overlay={
        \draw[
            fvdcyan,
            line width=0.5pt
        ]
        ([xshift=42mm,yshift=-13mm]frame.north west)
        --
        ([xshift=-7mm,yshift=-13mm]frame.north east);
    }
]
        \begin{itemize}
    \fi
}
{%
    \end{itemize}
    \ifdefined\HCode
        \HCode{%
</div>
</section>
        }%
    \else
        \end{tcolorbox}
    \fi
}


\newenvironment{leerdoelen}
{%
    \ifdefined\HCode
        \HCode{%
<section class="learning-box learning-box--goals">
<div class="learning-box__header">
<span class="learning-box__icon" aria-hidden="true"></span>
<span class="learning-box__title">Leerdoelen</span>
</div>
<div class="learning-box__content">
        }%
        \begin{itemize}
    \else
\begin{tcolorbox}[
    enhanced,
    breakable,
    colback=fvdmagentalight,
    colframe=fvdmagenta,
    coltext=fvdtext,
    boxrule=0.5pt,
    leftrule=4pt,
    arc=4mm,
    outer arc=4mm,
    left=7mm,
    right=6mm,
    top=5mm,
    bottom=5mm,
    before skip=6mm,
    after skip=6mm,
    title=Leerdoelen,
    coltitle=fvdmagenta,
    fonttitle=\bfseries\Large,
    attach title to upper={\par\smallskip},
    boxed title style={
        empty,
        left=0mm,
        right=0mm,
        top=0mm,
        bottom=0mm
    },
    overlay={
        \draw[
            fvdmagenta,
            line width=0.5pt
        ]
        ([xshift=42mm,yshift=-13mm]frame.north west)
        --
        ([xshift=-7mm,yshift=-13mm]frame.north east);
    }
]
        \begin{itemize}
    \fi
}
{%
    \end{itemize}
    \ifdefined\HCode
        \HCode{%
</div>
</section>
        }%
    \else
        \end{tcolorbox}
    \fi
}



\title{De regel van l'Hôpital}
\author{Ingmar Herreman}

\begin{document}

% FVD_CHAPTER_DATA_START
% Automatisch gegenereerd uit index.tex.
% Niet handmatig aanpassen.
\fvdchapterdata{3}{Afleiden van samengestelde functie}{3}
% FVD_CHAPTER_DATA_END

\begin{abstract}
Sommige limieten leveren na rechtstreekse substitutie geen getal op, maar
een onbepaalde vorm zoals \(0/0\) of \(\infty/\infty\).
De regel van l'Hôpital maakt het in zulke gevallen mogelijk om een
quotiënt te vervangen door een nieuw quotiënt van afgeleiden.

De belangrijkste vaardigheid is niet het afleiden zelf, maar het
\textbf{analyseren van de limiet}: welke vorm ontstaat, mag l'Hôpital
worden toegepast en is l'Hôpital wel de meest geschikte methode?
\end{abstract}

\maketitle

\begin{voorkennis}
  \item Je kan limieten algebraïsch bepalen.
  \item Je herkent de onbepaalde vormen \(0/0\) en \(\infty/\infty\).
  \item Je kan veeltermfuncties, rationale functies en irrationale functies afleiden.
  \item Je kan de productregel, quotiëntregel en kettingregel toepassen.
  \item Je kan algebraïsche uitdrukkingen ontbinden, vereenvoudigen en rationaliseren.
\end{voorkennis}

\begin{leerdoelen}
  \item Je kan bij een limiet eerst de vorm bepalen door substitutie.
  \item Je kan herkennen wanneer de regel van l'Hôpital rechtstreeks toepasbaar is.
  \item Je kan de regel van l'Hôpital correct toepassen bij de vormen \(0/0\) en \(\infty/\infty\).
  \item Je kan l'Hôpital meermaals toepassen wanneer opnieuw een geschikte onbepaalde vorm ontstaat.
  \item Je kan bij een onbepaalde vorm beslissen of l'Hôpital of een algebraïsche methode het meest geschikt is.
  \item Je kan een foutieve toepassing van l'Hôpital herkennen en verklaren.
  \item Je kan je keuze van een limietmethode wiskundig motiveren.
\end{leerdoelen}

% ============================================================
% 1. ONBEPAALDE VORMEN
% ============================================================

\FVDsection{Wanneer zegt substitutie nog niet genoeg?}

Bij het berekenen van een limiet is rechtstreekse substitutie meestal
de eerste stap. Soms levert die onmiddellijk de limietwaarde.

Bijvoorbeeld:

\[\lim_{x\to1}\frac{x+2}{x+3}=\frac34.\]

Hier is geen verdere bewerking nodig.

Maar beschouw nu

\[\lim_{x\to1}\frac{x^2-1}{x-1}.\]

Rechtstreekse substitutie geeft

\[\frac00.\]

Dit is geen limietwaarde. De teller en de noemer naderen beide nul, maar
daaruit kunnen we nog niet besluiten wat hun quotiënt doet.

\begin{definitie}[Onbepaalde vorm]
Een \textbf{onbepaalde vorm} is een symbolische vorm die op zichzelf
onvoldoende informatie bevat om de waarde van een limiet te bepalen.

Voor de regel van l'Hôpital zijn in dit hoofdstuk vooral

\[\boxed{\frac00 \qquad\text{en}\qquad \frac{\infty}{\infty}}\]

van belang.
\end{definitie}

Waarom noemen we \(0/0\) onbepaald? Vergelijk:

\[\lim_{x\to0}\frac{x}{x}=1,\]

\[\lim_{x\to0}\frac{x^2}{x}=0,\]

terwijl

\[\lim_{x\to0^+}\frac{x}{x^2}=+\infty.\]

In alle drie gevallen ontstaat bij substitutie formeel \(0/0\), maar de
limieten zijn verschillend. De vorm alleen bepaalt het antwoord dus niet.

\begin{opmerking}[De vorm is geen antwoord]
Schrijf nooit

\[\lim_{x\to a}f(x)=\frac00\]

alsof \(0/0\) de limietwaarde is.

Schrijf bijvoorbeeld:

\[
\frac{x^2-1}{x-1}\xrightarrow[x\to1]{}\frac00,
\]

dus de limiet heeft de onbepaalde vorm \(0/0\) en moet verder worden
onderzocht.
\end{opmerking}

% ============================================================
% 2. REGEL VAN L'HOPITAL
% ============================================================

\FVDsection{De regel van l'Hôpital}

De regel van l'Hôpital vergelijkt de verandering van teller en noemer
via hun afgeleiden.

\begin{eigenschap}[Regel van l'Hôpital]
Beschouw in een omgeving van \(a\) twee afleidbare functies \(f\) en
\(g\), met \(g'(x)\neq0\) voor \(x\neq a\).

Als

\[\lim_{x\to a}f(x)=\lim_{x\to a}g(x)=0\]

of als teller en noemer in grootte onbegrensd worden, en als

\[\lim_{x\to a}\frac{f'(x)}{g'(x)}=L\]

bestaat in \(\mathbb{R}\cup\{-\infty,+\infty\}\), dan geldt

\[\boxed{\lim_{x\to a}\frac{f(x)}{g(x)}=\lim_{x\to a}\frac{f'(x)}{g'(x)}=L.}\]

Dezelfde regel kan, onder de overeenkomstige voorwaarden, worden gebruikt
voor eenzijdige limieten en voor \(x\to+\infty\) of \(x\to-\infty\).
\end{eigenschap}

\begin{opmerking}[Geen quotiëntregel]
Bij l'Hôpital berekenen we \textbf{niet} de afgeleide van het quotiënt.

We vervangen

\[\frac{f(x)}{g(x)}\]

door

\[\frac{f'(x)}{g'(x)},\]

en dus niet door

\[\frac{f'(x)g(x)-f(x)g'(x)}{g(x)^2}.\]

De quotiëntregel en de regel van l'Hôpital zijn twee verschillende regels
met een verschillend doel.
\end{opmerking}

\FVDsubsection{Een betrouwbaar stappenplan}

Gebruik bij een mogelijke toepassing van l'Hôpital steeds dezelfde
denkvolgorde:

\begin{enumerate}
  \item \textbf{Substitueer} en bepaal de vorm.
  \item Controleer of de vorm \(0/0\) of \(\infty/\infty\) is.
  \item Alleen dan: \textbf{differentieer teller en noemer afzonderlijk}.
  \item Bepaal opnieuw de limiet.
  \item Ontstaat opnieuw \(0/0\) of \(\infty/\infty\)? Dan mag l'Hôpital opnieuw worden toegepast.
  \item Controleer of een algebraïsche methode niet eenvoudiger of inzichtelijker is.
\end{enumerate}

% ============================================================
% 3. VORM 0/0
% ============================================================

\FVDsection{De vorm \(\frac00\)}

\begin{voorbeeld}[Een eerste toepassing]
Bepaal

\[\lim_{x\to3}\frac{27-x^3}{2x^2-11x+15}.\]

Substitutie geeft

\[\frac{27-27}{18-33+15}=\frac00.\]

L'Hôpital is dus toepasbaar:

\[
\lim_{x\to3}\frac{27-x^3}{2x^2-11x+15}
\overset{\mathrm H}{=}
\lim_{x\to3}\frac{-3x^2}{4x-11}
=-27.
\]

Dus

\[\boxed{-27}.\]
\end{voorbeeld}

\begin{voorbeeld}[Met een wortelfunctie]
Bepaal

\[\lim_{x\to4}\frac{\sqrt{x}-2}{x-4}.\]

Substitutie geeft \(0/0\). Daarom mag l'Hôpital worden toegepast:

\[
\lim_{x\to4}\frac{\sqrt{x}-2}{x-4}
\overset{\mathrm H}{=}
\lim_{x\to4}\frac{\frac{1}{2\sqrt{x}}}{1}
=\frac14.
\]

Dus

\[\boxed{\frac14}.\]

Deze limiet kan ook worden berekend door teller en noemer met de
toegevoegde vorm te vermenigvuldigen. Dat is belangrijk: l'Hôpital is
een mogelijke methode, niet noodzakelijk de enige methode.
\end{voorbeeld}

\begin{voorbeeld}[Ontbinden of l'Hôpital?]
Bepaal

\[\lim_{x\to2}\frac{x^3-8}{x^2-4}.\]

Substitutie geeft \(0/0\), dus l'Hôpital mag:

\[
\lim_{x\to2}\frac{x^3-8}{x^2-4}
\overset{\mathrm H}{=}
\lim_{x\to2}\frac{3x^2}{2x}
=3.
\]

Maar ontbinden kan eveneens:

\[
\frac{x^3-8}{x^2-4}
=
\frac{(x-2)(x^2+2x+4)}{(x-2)(x+2)}
=
\frac{x^2+2x+4}{x+2}.
\]

Daaruit volgt opnieuw

\[\boxed{3}.\]

Beide methoden zijn correct. Hier zijn ze ongeveer even kort.
\end{voorbeeld}

% ============================================================
% 4. INFINITY/INFINITY
% ============================================================

\FVDsection{De vorm \(\frac{\infty}{\infty}\)}

De tweede vorm waarop l'Hôpital rechtstreeks kan worden toegepast, is
\(\infty/\infty\).

\begin{voorbeeld}[Een eenvoudige rationale functie]
Bepaal

\[\lim_{x\to+\infty}\frac{2x^2}{x+4}.\]

Teller en noemer worden beide onbegrensd. We hebben dus de vorm
\(\infty/\infty\).

Met l'Hôpital:

\[
\lim_{x\to+\infty}\frac{2x^2}{x+4}
\overset{\mathrm H}{=}
\lim_{x\to+\infty}4x
=+\infty.
\]

Dus

\[\boxed{+\infty}.\]
\end{voorbeeld}

\begin{voorbeeld}[Een limiet met waarde nul]
Bepaal

\[\lim_{x\to+\infty}\frac{3x^2-5x+1}{2x^3+x}.\]

De vorm is \(\infty/\infty\). L'Hôpital geeft

\[
\lim_{x\to+\infty}\frac{3x^2-5x+1}{2x^3+x}
\overset{\mathrm H}{=}
\lim_{x\to+\infty}\frac{6x-5}{6x^2+1}.
\]

Opnieuw ontstaat \(\infty/\infty\). We mogen l'Hôpital opnieuw toepassen:

\[
\lim_{x\to+\infty}\frac{6x-5}{6x^2+1}
\overset{\mathrm H}{=}
\lim_{x\to+\infty}\frac6{12x}
=0.
\]

Dus

\[\boxed{0}.\]

Voor deze rationale functie is delen door de hoogste macht sneller. Dit
voorbeeld is daarom vooral nuttig om de voorwaarden en het herhaald
toepassen van l'Hôpital te oefenen.
\end{voorbeeld}

% ============================================================
% 5. MEERMAALS
% ============================================================

\FVDsection{L'Hôpital meermaals toepassen}

Na één toepassing kan opnieuw \(0/0\) of \(\infty/\infty\) ontstaan.
Dan mag de regel opnieuw worden gebruikt.

\begin{voorbeeld}[Herhaald toepassen bij \(0/0\)]
Bepaal

\[\lim_{x\to1}\frac{x^3-3x+2}{(x-1)^2}.\]

Substitutie geeft \(0/0\). Een eerste toepassing levert

\[
\lim_{x\to1}\frac{x^3-3x+2}{(x-1)^2}
\overset{\mathrm H}{=}
\lim_{x\to1}\frac{3x^2-3}{2(x-1)}.
\]

Opnieuw krijgen we \(0/0\). Daarom:

\[
\lim_{x\to1}\frac{3x^2-3}{2(x-1)}
\overset{\mathrm H}{=}
\lim_{x\to1}\frac{6x}{2}
=3.
\]

Dus

\[\boxed{3}.\]
\end{voorbeeld}

\begin{voorbeeld}[Herhaald toepassen bij \(\infty/\infty\)]
Bepaal

\[\lim_{x\to+\infty}\frac{2x^3+x^2+1}{3x^4-2x}.\]

De vorm is \(\infty/\infty\). We krijgen:

\[
\begin{aligned}
\lim_{x\to+\infty}\frac{2x^3+x^2+1}{3x^4-2x}
&\overset{\mathrm H}{=}\lim_{x\to+\infty}\frac{6x^2+2x}{12x^3-2}\\
&\overset{\mathrm H}{=}\lim_{x\to+\infty}\frac{12x+2}{36x^2}\\
&\overset{\mathrm H}{=}\lim_{x\to+\infty}\frac{12}{72x}=0.
\end{aligned}
\]

Dus

\[\boxed{0}.\]

Ook hier is delen door de hoogste macht efficiënter. Een sterke
limietberekening bestaat dus niet uit zo vaak mogelijk l'Hôpital
toepassen, maar uit het kiezen van een geschikte methode.
\end{voorbeeld}

% ============================================================
% 6. WANNEER NIET?
% ============================================================

\FVDsection{Wanneer mag l'Hôpital niet?}

L'Hôpital is geen algemene rekenregel voor quotiënten. De vorm moet eerst
worden gecontroleerd.

\begin{voorbeeld}[Geen onbepaalde vorm]
Bepaal

\[\lim_{x\to0}\frac{x+1}{x+2}.\]

Substitutie geeft onmiddellijk

\[\frac12.\]

Er is geen onbepaalde vorm. L'Hôpital mag dus niet worden toegepast.

Een leerling die schrijft

\[
\lim_{x\to0}\frac{x+1}{x+2}
\overset{\mathrm H}{=}
\lim_{x\to0}\frac11=1
\]

past de regel foutief toe.
\end{voorbeeld}

\begin{voorbeeld}[De vorm \(\frac0{\text{niet-nul}}\)]
Bepaal

\[\lim_{x\to2}\frac{x-2}{x^2+1}.\]

Substitutie geeft

\[\frac05=0.\]

Dit is geen onbepaalde vorm. De limiet is onmiddellijk

\[\boxed{0}.\]
\end{voorbeeld}

\begin{voorbeeld}[De vorm \(\frac{\text{niet-nul}}0\)]
Onderzoek

\[\lim_{x\to2}\frac1{x-2}.\]

Hier ontstaat geen vorm \(0/0\). We moeten de eenzijdige limieten
onderzoeken:

\[\lim_{x\to2^-}\frac1{x-2}=-\infty,\]

\[\lim_{x\to2^+}\frac1{x-2}=+\infty.\]

De tweezijdige limiet bestaat dus niet. L'Hôpital is hier niet de juiste
methode.
\end{voorbeeld}

% ============================================================
% 7. EENZIJDIGE LIMIETEN
% ============================================================

\FVDsection{L'Hôpital bij eenzijdige limieten}

De regel kan ook bij eenzijdige limieten worden gebruikt, zolang de
voorwaarden aan de beschouwde zijde zijn vervuld.

\begin{voorbeeld}[Een wortelfunctie aan de rand van het domein]
Bepaal

\[\lim_{x\to0^+}\frac{\sqrt{x}}{x}.\]

De teller en noemer gaan beide naar nul. We hebben dus \(0/0\).

Voor \(x>0\) zijn de functies afleidbaar. Daarom:

\[
\lim_{x\to0^+}\frac{\sqrt{x}}{x}
\overset{\mathrm H}{=}
\lim_{x\to0^+}\frac{1}{2\sqrt{x}}
=+\infty.
\]

Dus

\[\boxed{+\infty}.\]

Merk op dat alleen de rechterlimiet zinvol is, omdat \(\sqrt{x}\) voor
reële \(x\) niet gedefinieerd is wanneer \(x<0\).
\end{voorbeeld}

\begin{voorbeeld}[Het teken blijft belangrijk]
Bepaal

\[\lim_{x\to1^-}\frac{x^2-1}{(x-1)^2}.\]

Substitutie geeft \(0/0\). L'Hôpital geeft

\[
\lim_{x\to1^-}\frac{x^2-1}{(x-1)^2}
\overset{\mathrm H}{=}
\lim_{x\to1^-}\frac{2x}{2(x-1)}
=
\lim_{x\to1^-}\frac{x}{x-1}.
\]

Voor \(x\to1^-\) is de teller positief en nadert de noemer \(0\) langs
negatieve waarden. Daarom

\[\boxed{-\infty}.\]

De richting van de limiet mag dus niet verloren gaan na het toepassen van
l'Hôpital.
\end{voorbeeld}

% ============================================================
% 8. ANDERE VORMEN
% ============================================================

\FVDsection{Andere onbepaalde vormen: eerst algebraïsch omvormen}

Niet elke onbepaalde vorm is rechtstreeks geschikt voor l'Hôpital.
Binnen de functietypes die we nu kennen, is vooral de vorm
\(\infty-\infty\) belangrijk.

\begin{voorbeeld}[De vorm \(\infty-\infty\)]
Bepaal

\[\lim_{x\to+\infty}\left(\sqrt{x^2+x}-x\right).\]

Rechtstreeks ontstaat

\[\infty-\infty.\]

L'Hôpital kan niet rechtstreeks op een verschil worden toegepast.
We gebruiken eerst de toegevoegde vorm:

\[
\sqrt{x^2+x}-x
=
\frac{(\sqrt{x^2+x}-x)(\sqrt{x^2+x}+x)}
{\sqrt{x^2+x}+x}
=
\frac{x}{\sqrt{x^2+x}+x}.
\]

Nu ontstaat een quotiënt van de vorm \(\infty/\infty\). We kunnen
l'Hôpital toepassen, maar algebraïsch vereenvoudigen is hier sneller:

\[
\frac{x}{\sqrt{x^2+x}+x}
=
\frac{1}{\sqrt{1+\frac1x}+1}.
\]

Dus

\[\boxed{\frac12}.\]
\end{voorbeeld}

\begin{opmerking}[Eerst omvormen, dan beslissen]
Bij een vorm zoals \(\infty-\infty\) is de denkstap belangrijker dan de
regel zelf:

\[
\text{onbepaalde vorm}
\longrightarrow
\text{algebraïsch omvormen}
\longrightarrow
\text{nieuwe vorm onderzoeken}
\longrightarrow
\text{geschikte methode kiezen}.
\]
\end{opmerking}

% ============================================================
% 9. METHODEKEUZE
% ============================================================

\FVDsection{L'Hôpital of een andere methode?}

Dat l'Hôpital \emph{mag} worden toegepast, betekent niet automatisch dat
l'Hôpital de beste methode is.

\begin{opmerking}[Kies bewust]
Vraag jezelf achtereenvolgens af:

\begin{enumerate}
  \item Kan ik rechtstreeks substitueren?
  \item Ontstaat een onbepaalde vorm?
  \item Kan ik eenvoudig ontbinden, vereenvoudigen, rationaliseren of delen door de hoogste macht?
  \item Is de vorm \(0/0\) of \(\infty/\infty\), zodat l'Hôpital rechtstreeks mag?
  \item Welke methode is het kortst en maakt de structuur van de limiet het duidelijkst?
\end{enumerate}
\end{opmerking}

\begin{voorbeeld}[Ontbinden is hier natuurlijker]
Bepaal

\[\lim_{x\to2}\frac{x^2-4}{x-2}.\]

Met ontbinden:

\[
\frac{x^2-4}{x-2}
=
\frac{(x-2)(x+2)}{x-2}
=x+2
\qquad(x\neq2).
\]

Dus

\[\boxed{4}.\]

L'Hôpital is eveneens toegestaan:

\[
\lim_{x\to2}\frac{x^2-4}{x-2}
\overset{\mathrm H}{=}
\lim_{x\to2}2x
=4.
\]

Ontbinden toont hier echter onmiddellijk waarom de schijnbare
singulariteit verdwijnt.
\end{voorbeeld}

\begin{voorbeeld}[Delen door de hoogste macht is hier natuurlijker]
Bepaal

\[\lim_{x\to+\infty}\frac{3x^2+2x-1}{x^2+5}.\]

De vorm is \(\infty/\infty\), dus l'Hôpital mag. Maar delen door \(x^2\)
geeft onmiddellijk

\[
\lim_{x\to+\infty}
\frac{3+\frac2x-\frac1{x^2}}
{1+\frac5{x^2}}
=3.
\]

Dus

\[\boxed{3}.\]

De algebraïsche methode is hier korter en laat meteen zien dat de
coëfficiënten van de hoogste machten bepalend zijn.
\end{voorbeeld}

\begin{voorbeeld}[L'Hôpital is hier bijzonder efficiënt]
Bepaal

\[\lim_{x\to4}\frac{\sqrt{x}-2}{x-4}.\]

Rationaliseren is mogelijk, maar l'Hôpital geeft onmiddellijk

\[
\lim_{x\to4}\frac{\sqrt{x}-2}{x-4}
\overset{\mathrm H}{=}
\lim_{x\to4}\frac1{2\sqrt{x}}
=\frac14.
\]

Hier is l'Hôpital dus een zeer efficiënte keuze.
\end{voorbeeld}

% ============================================================
% 10. KORTE CONTROLE
% ============================================================

\FVDsection{De eerste oefeningen}

\begin{oefening}[Mag l'Hôpital?]{\oefB\ESS}

Bepaal bij elke limiet eerst de vorm. Beslis daarna of l'Hôpital
\textbf{rechtstreeks} mag worden toegepast. Motiveer kort.

\begin{question}
\[\lim_{x\to1}\frac{x^2-1}{x-1}.\]

\begin{oplossing}
Substitutie geeft \(0/0\). L'Hôpital mag rechtstreeks worden toegepast.
\end{oplossing}
\end{question}

\begin{question}
\[\lim_{x\to0}\frac{x+1}{x+2}.\]

\begin{oplossing}
Substitutie geeft \(1/2\). Dit is geen onbepaalde vorm. L'Hôpital mag
niet worden toegepast.
\end{oplossing}
\end{question}

\begin{question}
\[\lim_{x\to+\infty}\frac{2x^2+1}{x^3-4}.\]

\begin{oplossing}
Teller en noemer worden in grootte onbegrensd. De vorm is
\(\infty/\infty\), dus l'Hôpital mag rechtstreeks worden toegepast.
\end{oplossing}
\end{question}

\begin{question}
\[\lim_{x\to+\infty}\left(\sqrt{x^2+x}-x\right).\]

\begin{oplossing}
De vorm is \(\infty-\infty\). L'Hôpital mag niet rechtstreeks worden
toegepast. De uitdrukking moet eerst worden omgevormd, bijvoorbeeld met
de toegevoegde vorm.
\end{oplossing}
\end{question}

\begin{question}
\[\lim_{x\to2}\frac1{x-2}.\]

\begin{oplossing}
Dit is geen vorm \(0/0\) of \(\infty/\infty\). L'Hôpital mag niet worden
toegepast. De eenzijdige limieten moeten afzonderlijk worden onderzocht.
\end{oplossing}
\end{question}
\end{oefening}

\begin{oefening}[Bereken en verantwoord]{\oefU\ESS}
  
Bereken de limieten. Noteer telkens eerst de vorm en motiveer je
methodekeuze.

\begin{question}
\[\lim_{x\to2}\frac{x^3-8}{x^2-4}.\]

\begin{oplossing}
De vorm is \(0/0\). Zowel ontbinden als l'Hôpital is geschikt.

Met l'Hôpital:

\[
\lim_{x\to2}\frac{x^3-8}{x^2-4}
\overset{\mathrm H}{=}
\lim_{x\to2}\frac{3x^2}{2x}
=3.
\]

Dus

\[\boxed{3}.\]
\end{oplossing}
\end{question}

\begin{question}
\[\lim_{x\to4}\frac{x-4}{\sqrt{x}-2}.\]

\begin{oplossing}
De vorm is \(0/0\). Met l'Hôpital:

\[
\lim_{x\to4}\frac{x-4}{\sqrt{x}-2}
\overset{\mathrm H}{=}
\lim_{x\to4}\frac{1}{1/(2\sqrt{x})}
=4.
\]

Dus

\[\boxed{4}.\]

Ook algebraïsch vereenvoudigen met
\(x-4=(\sqrt{x}-2)(\sqrt{x}+2)\) is zeer efficiënt.
\end{oplossing}
\end{question}

\begin{question}
\[\lim_{x\to+\infty}\frac{5x^2-x}{2x^2+3}.\]

\begin{oplossing}
De vorm is \(\infty/\infty\). L'Hôpital mag, maar delen door \(x^2\)
is hier eenvoudiger:

\[
\lim_{x\to+\infty}\frac{5-\frac1x}{2+\frac3{x^2}}
=\frac52.
\]

Dus

\[\boxed{\frac52}.\]
\end{oplossing}
\end{question}

\begin{question}
Een leerling gebruikt l'Hôpital voor

\[\lim_{x\to2}\frac{x^2+1}{x+1}.\]

Analyseer de aanpak.

\begin{oplossing}
Substitutie geeft

\[\frac53.\]

Dit is geen onbepaalde vorm. De limiet is dus onmiddellijk
\(\frac53\) en l'Hôpital mag niet worden toegepast.
\end{oplossing}
\end{question}
\end{oefening}

\begin{oefening}[Analyseer de methode]{\oefU\ESS}

Drie leerlingen berekenen

\[\lim_{x\to+\infty}\frac{4x^2+x}{2x^2-3}.\]

Leerling A past l'Hôpital toe. Leerling B deelt teller en noemer door
\(x^2\). Leerling C zegt onmiddellijk dat de limiet \(4/2=2\) is omdat
teller en noemer dezelfde graad hebben.

\begin{question}
Welke methoden zijn wiskundig geldig? Motiveer.

\begin{oplossing}
Alle drie kunnen tot de correcte limiet leiden.

Leerling A mag l'Hôpital toepassen omdat de vorm
\(\infty/\infty\) is:

\[
\lim_{x\to+\infty}\frac{4x^2+x}{2x^2-3}
\overset{\mathrm H}{=}
\lim_{x\to+\infty}\frac{8x+1}{4x}
=2.
\]

Leerling B gebruikt de standaardmethode voor rationale functies:

\[
\frac{4x^2+x}{2x^2-3}
=
\frac{4+\frac1x}{2-\frac3{x^2}}
\longrightarrow2.
\]

Leerling C gebruikt een gekende eigenschap van rationale functies met
gelijke graad. Die conclusie is correct als de eigenschap eerder
voldoende werd verantwoord.

Voor deze limiet is de aanpak van leerling B het meest transparant:
l'Hôpital is toegestaan, maar niet nodig.
\end{oplossing}
\end{question}
\end{oefening}



\end{document}
